\documentclass[12pt,letterpaper]{article}

\usepackage[margin=1in]{geometry}
\usepackage[T1]{fontenc}
\usepackage[utf8]{inputenc}
\usepackage{lmodern}
\usepackage{microtype}
\usepackage{setspace}
\usepackage{enumitem}
\usepackage{titlesec}
\usepackage{xcolor}
\usepackage[round,authoryear]{natbib}
\usepackage[hidelinks]{hyperref}

\onehalfspacing

\setlist[itemize]{leftmargin=*,noitemsep,topsep=1pt}
\setlist[enumerate]{leftmargin=*,noitemsep,topsep=1pt}
\titleformat{\section}{\large\bfseries}{\thesection}{0pt}{}
\titlespacing*{\section}{0pt}{0.5ex}{0.2ex}
\setlength{\bibsep}{0pt}
\urlstyle{same}

\newcommand{\placeholder}[1]{\textcolor{gray}{[\textit{#1}]}}

\begin{document}

\begin{center}
{\Large\bfseries The Effect of Feedback Modality on Speed--Accuracy Tradeoff in Target-Selection Tasks\par}
\vspace{0.15em}
{\normalsize MIE~286 (Winter~2026) --- Term Project Proposal\par}
\vspace{0.3em}
{\normalsize
\textbf{Team:} \placeholder{Team \#} \quad
\textbf{Members:} \placeholder{Name 1}, \placeholder{Name 2}, \placeholder{Name 3}\par}
{\normalsize \today\par}
\end{center}
\vspace{-0.2em}
\hrule
\vspace{0.4em}

\section*{1. Project Overview}
Rapid target-selection (pointing and clicking) tasks exhibit a speed--accuracy tradeoff and are commonly modeled with Fitts-style formulations \citep{fitts1954,mackenzie1992}. The tradeoff can be interpreted as people adjusting internal response criteria (``go faster'' vs.\ ``be more accurate'') \citep{heitz2014,ratcliff2008}. Theory and prior evidence suggest that presenting secondary information in the same modality as a primary task can increase interference, while distributing information across modalities can reduce competition for attention and working memory \citep{wickens2015,klingner_nd}. Auditory feedback can communicate performance without requiring additional visual sampling, and is widely used in augmented feedback/sonification contexts \citep{sigrist2013,hermann2011}. We will test whether feedback modality shifts the speed--accuracy tradeoff in a rapid clicking task.

\noindent\textbf{Design summary:} IV = Feedback Type (visual vs.\ auditory); DVs = response time (ms) and radial error (px); within-subject with counterbalanced order; $N\ge 24$ adults (18+).

\section*{2. Research Question and Hypotheses}
\textbf{Research question:} Does real-time performance feedback modality (visual vs.\ auditory) influence the speed--accuracy tradeoff in a visual target-selection task?
\begin{itemize}
  \item \textbf{H1 (Speed):} Mean response time will be lower under auditory feedback than under visual feedback \citep{wickens2015}.
  \item \textbf{H2 (Accuracy):} Auditory feedback will increase speed but slightly worsen accuracy (higher radial error), consistent with a classic speed--accuracy tradeoff \citep{heitz2014}.
  \item \textbf{H3 (Tradeoff structure):} The RT--error relationship (correlation/slope) will differ by condition \citep{heitz2014}.
\end{itemize}

\section*{3. Methods}
\textbf{Participants.} We will recruit at least 24 participants (18+), with normal/corrected vision and no self-reported hearing impairments. We will record age range and device type (mouse/trackpad).

\textbf{Task and apparatus.} Participants complete a browser-based target-selection task on a fixed-size canvas (e.g., $800\times 500$ px). Each trial shows one circular target (fixed radius; e.g., 40 px) at a random location; participants click as quickly and accurately as possible \citep{mackenzie1992}. Click timestamps and coordinates are logged; the next target appears after a fixed inter-trial interval (e.g., 500 ms).

\textbf{Feedback manipulation (two active conditions).}
\begin{itemize}
  \item \textbf{Visual feedback:} Color flash after each click (green hit / red miss) plus a small HUD showing running mean RT, hit rate, and mean radial error.
  \item \textbf{Auditory feedback:} No HUD (trial counter only) plus a short tone after each click: high pitch fast hit, medium pitch slow hit, low ``thud'' miss; the fast/slow cutoff will be set from a short pilot \citep{hermann2011}.
\end{itemize}

\textbf{Procedure.} Consent and demographics $\rightarrow$ instructions $\rightarrow$ 10 practice trials $\rightarrow$ Block~1 (60 trials) $\rightarrow$ short break $\rightarrow$ Block~2 (60 trials). Condition order is counterbalanced by random assignment to Visual$\rightarrow$Auditory vs Auditory$\rightarrow$Visual.

\textbf{Measures.} Response time (ms) is measured from target onset to click. Radial error (px) is the Euclidean distance from click to target center (hit if error $\le$ radius).

\section*{4. Analysis Plan (R/RStudio)}
\begin{enumerate}
  \item \textbf{Cleaning:} exclude implausible RTs (e.g., $<$200 ms or $>$5000 ms); verify expected trial counts.
  \item \textbf{Summary:} per participant$\times$condition compute mean RT, mean radial error, and hit rate.
  \item \textbf{H1--H2:} paired $t$-tests (or Wilcoxon signed-rank if non-normal) comparing visual vs.\ auditory for mean RT and mean radial error.
  \item \textbf{H3:} quantify RT--error association by condition (correlation) and test a moderation model \texttt{mean\_error} $\sim$ \texttt{mean\_rt} $\times$ \texttt{condition}.
\end{enumerate}

\begin{thebibliography}{9}\setlength{\itemsep}{0pt}
\bibitem[Fitts(1954)]{fitts1954}
Fitts, P.~M. (1954). The information capacity of the human motor system in controlling the amplitude of movement. \emph{Journal of Experimental Psychology}, 47(6), 381--391.

\bibitem[MacKenzie(1992)]{mackenzie1992}
MacKenzie, I.~S. (1992). Fitts' law as a research and design tool in human--computer interaction. \emph{Human--Computer Interaction}, 7(1), 91--139.

\bibitem[Heitz(2014)]{heitz2014}
Heitz, R.~P. (2014). The speed--accuracy tradeoff: history, physiology, methodology, and behavior. \emph{Frontiers in Neuroscience}, 8, 150.

\bibitem[Ratcliff and McKoon(2008)]{ratcliff2008}
Ratcliff, R., \& McKoon, G. (2008). The diffusion decision model: theory and data for two-choice decision tasks. \emph{Neural Computation}, 20(4), 873--922.

\bibitem[Wickens et~al.(2015)]{wickens2015}
Wickens, C.~D., Hollands, J.~G., Banbury, S., \& Parasuraman, R. (2015). \emph{Engineering Psychology and Human Performance} (4th ed.). Routledge.

\bibitem[Sigrist et~al.(2013)]{sigrist2013}
Sigrist, R., Rauter, G., Riener, R., \& Wolf, P. (2013). Augmented visual, auditory, haptic feedback in motor learning: A review. \emph{Psychonomic Bulletin \& Review}, 20(1), 21--53.

\bibitem[Hermann et~al.(2011)]{hermann2011}
Hermann, T., Hunt, A., \& Neuhoff, J.~G. (Eds.). (2011). \emph{The Sonification Handbook}. Logos Verlag.

\bibitem[Klingner et~al.(n.d.)]{klingner_nd}
Klingner, J., Tversky, B., \& Hanrahan, P. (n.d.). \emph{Aural vs.\ Visual Cognitive Load}. Available at \url{https://graphics.stanford.edu/papers/visual-cognitive-load/klingner_aural_vs_visual_cognitive_load.pdf} (accessed 2026-02-15).
\end{thebibliography}

\end{document}
